
\setcounter{footnote}{73}

\section*{\SGAInlineRefTarget{sga3:iii:section:4}4. Infinitesimal extensions of closed subgroups}
\addcontentsline{toc}{section}{4. Infinitesimal extensions of closed subgroups}

We first state a result valid in an arbitrary abelian category.

\medskip
\noindent\SGAInlineRefTarget{sga3:iii:lemma:4-1}\textbf{Lemma 4.1.}
\textit{Let
\[
 0\longrightarrow A'\xrightarrow{i}A\xrightarrow{p}A''\longrightarrow0
\]
be an exact sequence, let \(\phi:A'\to Q\) be a morphism, and let
\(\pi:A''\to P\) be an epimorphism with kernel \(C\). Let \(E\) be the
set, up to isomorphism, of quadruples \((B,f,g,h)\) such that the sequence
\[
 0\longrightarrow Q\xrightarrow{f}B\xrightarrow{g}P\longrightarrow0
\]
is exact and the following diagram commutes:}

\begin{center}
  \SGAThreeNativeInput[0.78\textwidth]
    {native_diagrams/exp3/diagram_031_exp3_p042_lemma41_commutative.tex}
\end{center}

\begin{enumerate}
\item[(i)] \textit{The set \(E\) is nonempty if and only if the image in
  \(\operatorname{Ext}^1(C,Q)\) of the element \(A\) of
  \(\operatorname{Ext}^1(A'',A')\) is zero.}
\item[\SGAInlineRefTarget{sga3:iii:lemma:4-1:item:ii}(ii)] \textit{Under these conditions, \(E\) is a principal homogeneous
  set under the abelian group \(\operatorname{Hom}(C,Q)\).}
\end{enumerate}

Introduce the amalgamated sum
\[
 B'=A\mathop{\sqcup}\limits_{A'}Q.
\]
There is then a commutative diagram with exact rows:
\footnote{Editor's note: One has
\(\operatorname{Coker}(j)=B'\mathop{\sqcup}\limits_Q0
=A\mathop{\sqcup}\limits_{A'}0=A''\), and
\(\operatorname{Ker}(j)\simeq\operatorname{Ker}(i)=0\), as one sees by
reasoning ``as if the category were a category of modules.'' For a proof
entirely in terms of arrows, see for example [Fr64], Theorem 2.5.4 (*).}

\begin{center}
  \SGAThreeNativeInput[0.78\textwidth]
    {native_diagrams/exp3/diagram_032_exp3_p043_lemma41_pushout.tex}
\end{center}

It is clear that the category of solutions to the problem is canonically
isomorphic to the category of solutions to the corresponding problem for
the sequence
\[
 0\longrightarrow Q\longrightarrow B'\longrightarrow A''\longrightarrow0
\]
and the morphisms \(\operatorname{id}_Q\) and \(\pi:A''\to P\).
\footnote{Editor's note: In what follows, \(A\) has been replaced by
\(B'\), and the end of the argument has been expanded.}
In this case, \(E\) is in bijection with the set of subobjects \(N\) of
\(B'\) for which \(B'\to A''\) induces an isomorphism of \(N\) with the
kernel \(C\) of \(A''\to P\); equivalently, it is the set of morphisms
\(e:C\to B'\) lifting the canonical morphism \(C\to A''\). The abelian
group \(G=\operatorname{Hom}(C,Q)\) acts on \(E\) by
\(g\cdot e=g+e\), with addition in \(\operatorname{Hom}(C,B')\); if
\(E\ne\varnothing\), this makes \(E\) a principal homogeneous set under
\(G\).

We deduce the following.

\medskip
\noindent\SGAInlineRefTarget{sga3:iii:proposition:4-2}\textbf{Proposition 4.2.}
\footnote{Editor's note: The statement has been rewritten so as to fit
exactly the application made in \SGARefLink{sga3:iii:proposition:4-3}{4.3}; the proof has also been expanded,
following indications supplied by M. Demazure.}
\textit{Let \(S\) be a scheme, let \(S_{\mathscr J}\) be the closed
subscheme defined by a quasi-coherent square-zero ideal \(\mathscr J\),
let \(X\) be an \(S\)-scheme, and let \(\mathscr F\) be an
\(\mathscr O_X\)-module. Put
\[
 X_{\mathscr J}=X\times_S S_{\mathscr J},\qquad
 \mathscr F_{\mathscr J}
 =\mathscr F\otimes_{\mathscr O_S}\mathscr O_{S_{\mathscr J}},
\]
and let
\(\mathscr G_{\mathscr J}=\mathscr F_{\mathscr J}/\mathscr H_{\mathscr J}\)
be a quotient module of \(\mathscr F_{\mathscr J}\). Suppose given a
morphism of \(\mathscr O_{X_{\mathscr J}}\)-modules
\[
 f:\mathscr J\otimes_{\mathscr O_{S_{\mathscr J}}}
       \mathscr G_{\mathscr J}\longrightarrow\mathscr Q.
\]
Let \(\mathcal E\) be the sheaf of sets on \(X\) defined as follows: for
each open subset \(U\) of \(X\), \(\mathcal E(U)\) is the set of quotient
modules \(\mathscr G\) of \(\mathscr F|_U\) such that
\(\mathscr G/\mathscr J\mathscr G=\mathscr G_{\mathscr J}|_U\) and for
which there exists an isomorphism
\[
 h:\mathscr J\mathscr G\xrightarrow{\sim}\mathscr Q|_U
\]
making the following diagram commute:}

\begin{center}
  \SGAThreeNativeInput[0.43\textwidth]
    {native_diagrams/exp3/diagram_033_exp3_p043_prop42_triangle.tex}
\end{center}

\textit{The isomorphism \(h\) is then unique, since
\(
 \mathscr J\otimes_{\mathscr O_{S_{\mathscr J}}}
 (\mathscr G_{\mathscr J}|_U)\to\mathscr J\mathscr G
\)
is an epimorphism. Then \(\mathcal E\) is a sheaf formally principal
homogeneous under the sheaf of commutative groups
\[
 \mathscr A=
 \operatorname{Hom}_{\mathscr O_X}(\mathscr H_{\mathscr J},\mathscr Q)
 =\operatorname{Hom}_{\mathscr O_{X_{\mathscr J}}}
   (\mathscr H_{\mathscr J},\mathscr Q).
\]}

\noindent\textit{Proof.}
If \(\mathcal E(U)=\varnothing\), there is nothing to prove; we may
therefore suppose that \(\mathcal E(U)\) contains an element
\(\widetilde{\mathscr G}\). In the following diagram, \(h\) is an
isomorphism and all the arrows are epimorphisms:

\begin{center}
  \SGAThreeNativeInput[0.70\textwidth]
    {native_diagrams/exp3/diagram_034_exp3_p044_prop42_epimorphisms.tex}
\end{center}

Thus the morphism
\[
 \mathscr J\otimes_{\mathscr O_{S_{\mathscr J}}}
 (\mathscr F_{\mathscr J}|_U)\longrightarrow\mathscr Q|_U
\]
induces a necessarily unique epimorphism
\(\phi:\mathscr J\mathscr F|_U\to\mathscr Q|_U\). If \(\mathscr G\)
is an \(\mathscr O_U\)-module such that
\(\mathscr G/\mathscr J\mathscr G=\mathscr G_{\mathscr J}|_U\) and there
is a commutative diagram with exact rows

\begin{center}
  \SGAThreeNativeInput[0.69\textwidth]
    {native_diagrams/exp3/diagram_035_exp3_p044_prop42_exact_rows.tex}
\end{center}

\noindent(where \(p_{\mathscr J}\) is the projection
\(\mathscr G\to\mathscr G/\mathscr J\mathscr G
=\mathscr G_{\mathscr J}|_U\), so that
\(\mathscr Q|_U=\operatorname{Ker}(p_{\mathscr J})
=\mathscr J\mathscr G\)), then \(\mathscr G\) may be identified with a
quotient module of \(\mathscr F|_U\). Consequently, by \SGARefLink{sga3:iii:lemma:4-1:item:ii}{4.1(ii)}, the set
\(\mathcal E(U)\) is principal homogeneous under the abelian group
\[
 \operatorname{Hom}_{\mathscr O_X}
  (\mathscr H_{\mathscr J},\mathscr Q)(U)
 =\operatorname{Hom}_{\mathscr O_{X_{\mathscr J}}}
  (\mathscr H_{\mathscr J},\mathscr Q)(U).
\]

\medskip
\noindent\SGAInlineRefTarget{sga3:iii:proposition:4-3}\textbf{Proposition 4.3.} \textit{(TDTE \SGARefLink{sga3:iv:subsection:5:1}{IV 5.1}) Let \(S\) be a
scheme, let \(S_{\mathscr J}\) be the closed subscheme defined by a
quasi-coherent square-zero ideal \(\mathscr J\), let \(X\) be an
\(S\)-scheme, and let \(\mathscr F\) be a quasi-coherent
\(\mathscr O_X\)-module. Put
\[
 X_{\mathscr J}=X\times_S S_{\mathscr J},\qquad
 \mathscr F_{\mathscr J}
 =\mathscr F\otimes_{\mathscr O_S}\mathscr O_{S_{\mathscr J}}.
\]
Let
\(\mathscr G_{\mathscr J}=\mathscr F_{\mathscr J}/\mathscr H_{\mathscr J}\)
be a quasi-coherent quotient module of \(\mathscr F_{\mathscr J}\), flat
over \(S_{\mathscr J}\). For each open subset \(U\) of \(X\), let
\(\mathcal E(U)\) be the set of quasi-coherent quotient modules
\(\mathscr G\)\footnote{Editor's note: To lighten the statement, the
hypothesis that \(\mathscr G\) be quasi-coherent has been added here; the
remark that this hypothesis is automatically satisfied has been moved to
the proof, which has been expanded accordingly.} of \(\mathscr F|_U\),
flat over \(S\), such that
\(\mathscr G/\mathscr J\mathscr G\simeq
\mathscr G_{\mathscr J}|_U\). Then the \(\mathcal E(U)\) form a sheaf of
sets \(\mathcal E\) on \(X\), formally principal homogeneous under the
sheaf of commutative groups
\[
 \mathscr A=
 \operatorname{Hom}_{\mathscr O_{X_{\mathscr J}}}
  \left(
   \mathscr H_{\mathscr J},
   \mathscr J\otimes_{\mathscr O_{S_{\mathscr J}}}
   \mathscr G_{\mathscr J}
  \right).
\]}

\noindent\textit{Proof.}
Let \(\pi:X\to S\) be the structural morphism. Let \(U\) be an open
subset of \(X\), and let \(\mathscr G\) be an \(\mathscr O_U\)-module
flat over \(S\), such that
\(\mathscr G/\mathscr J\mathscr G\simeq
\mathscr G_{\mathscr J}|_U\). For every \(x\in U\), the stalk
\(\mathscr G_x\) is a flat module over the local ring
\(\mathscr O_{S,s}\), where \(s=\pi(x)\). Hence the morphism
\[
 \mathscr J_s\otimes_{\mathscr O_{S,s}}
  (\mathscr G/\mathscr J\mathscr G)_x
 =\mathscr J_s\otimes_{\mathscr O_{S,s}}\mathscr G_x
 \longrightarrow(\mathscr J\mathscr G)_x
\]
is bijective. We therefore have an exact sequence
\[
 0\longrightarrow
 \mathscr J\otimes_{\mathscr O_S}(\mathscr G_{\mathscr J}|_U)
 \longrightarrow\mathscr G\longrightarrow
 \mathscr G_{\mathscr J}|_U\longrightarrow0.
\]

Since
\(\mathscr J\otimes_{\mathscr O_S}(\mathscr G_{\mathscr J}|_U)\) and
\(\mathscr G_{\mathscr J}|_U\) are quasi-coherent
\(\mathscr O_U\)-modules, \(\mathscr G\) is quasi-coherent as well
(EGA III, 1.4.17).

Conversely, since \(\mathscr G_{\mathscr J}\) is flat over
\(S_{\mathscr J}\), if \(\mathscr G\) is a quasi-coherent
\(\mathscr O_U\)-module such that
\(\mathscr G/\mathscr J\mathscr G\simeq\mathscr G_{\mathscr J}\) and
the morphism
\[
 \mathscr J\otimes_{\mathscr O_{S_{\mathscr J}}}
 \mathscr G_{\mathscr J}\longrightarrow\mathscr J\mathscr G
\]
is bijective, then \(\mathscr G\) is flat over \(S\), by the
``fundamental criterion of flatness'' (SGA 1, IV, 5.5).
\footnote{Editor's note: See also [BAC], \S III.5, Theorem 1.}

Consequently, the set \(\mathcal E(U)\) considered here is the set
considered in \SGARefLink{sga3:iii:proposition:4-2}{4.2} when \(f\) is the identity morphism of
\(\mathscr J\otimes_{\mathscr O_{S_{\mathscr J}}}
\mathscr G_{\mathscr J}\); the conclusion therefore follows from \SGARefLink{sga3:iii:proposition:4-2}{4.2}.
\hfill\(\square\)

Retain the preceding notation.
\footnote{Editor's note: The discussion that follows has been expanded and
\SGARefLink{sga3:iii:corollary:4-3-1}{Corollary 4.3.1} has been added. Recall also that ``pseudo-torsor'' is
synonymous with ``formally principal homogeneous.''}
Let \(Y_{\mathscr J}\) be a closed subscheme of \(X_{\mathscr J}\),
defined by a quasi-coherent ideal \(\mathscr I_{Y_{\mathscr J}}\), and
suppose that \(Y_{\mathscr J}\) is flat over \(S_{\mathscr J}\). Applying
\SGARefLink{sga3:iii:proposition:4-3}{4.3} to \(\mathscr F=\mathscr O_X\) and
\[
 \mathscr G_{\mathscr J}=\mathscr O_{Y_{\mathscr J}}
 =\mathscr O_{X_{\mathscr J}}/\mathscr I_{Y_{\mathscr J}},
\]
we obtain the following corollary.

\medskip
\noindent\SGAInlineRefTarget{sga3:iii:corollary:4-3-1}\textbf{Corollary 4.3.1.}
\textit{Let \(S,S_{\mathscr J},\mathscr J,X,X_{\mathscr J},
Y_{\mathscr J}\), and \(\mathscr I_{Y_{\mathscr J}}\) be as above, and
suppose that \(Y_{\mathscr J}\) is flat over \(S_{\mathscr J}\). Let
\(\mathscr A_{\mathscr J}\) be the sheaf of commutative groups
\[
 \operatorname{Hom}_{\mathscr O_{X_{\mathscr J}}}
 \left(
  \mathscr I_{Y_{\mathscr J}},
  \mathscr J\otimes_{\mathscr O_{S_{\mathscr J}}}
  \mathscr O_{Y_{\mathscr J}}
 \right)
\]
on \(X_{\mathscr J}\), and put
\(\mathscr A=i_*(\mathscr A_{\mathscr J})\), where
\(i:X_{\mathscr J}\hookrightarrow X\) is the immersion. For each open
subset \(U\) of \(X\), let \(\mathcal E(U)\) be the set of closed
subschemes \(Y\) of \(U\), flat over \(S\), such that
\[
 Y\times_S S_{\mathscr J}=Y_{\mathscr J}\cap U.
\]
Then \(\mathcal E\) is an \(\mathscr A\)-pseudo-torsor.}

\textit{If, moreover, such a \(Y\) exists locally---that is, if every
\(x\in X\) has an open neighborhood \(U\) with
\(\mathcal E(U)\ne\varnothing\)---then \(\mathcal E\) is an
\(\mathscr A\)-torsor.} The \(\mathscr A\)-torsors on \(X\) are
parametrized by
\[
 H^1(X,\mathscr A)=H^1(X_{\mathscr J},\mathscr A_{\mathscr J})
\]
(see, for example, EGA IV\textsubscript{4}, 16.5.15), and
\(\mathcal E\) has a global section---that is,
\(\mathcal E(X)\ne\varnothing\)---if and only if its corresponding
cohomology class is zero. Thus we obtain:

\medskip
\noindent\SGAInlineRefTarget{sga3:iii:corollary:4-4}\textbf{Corollary 4.4.}
\textit{Let \(S,S_{\mathscr J},\mathscr J,X,X_{\mathscr J},
Y_{\mathscr J}\), and \(\mathscr I_{Y_{\mathscr J}}\) be as above, and
suppose that \(Y_{\mathscr J}\) is flat over \(S_{\mathscr J}\). Let
\(E\) be the set of closed subschemes \(Y\) of \(X\), flat over \(S\),
such that \(Y\times_S S_{\mathscr J}=Y_{\mathscr J}\).}
\begin{enumerate}
\item[(i)] \textit{The set \(E\) is either empty or principal homogeneous
under the abelian group}
\[
 H^0(X,\mathscr A)=H^0(X_{\mathscr J},\mathscr A_{\mathscr J})
 =\operatorname{Hom}_{\mathscr O_{X_{\mathscr J}}}
 \left(
  \mathscr I_{Y_{\mathscr J}},
  \mathscr J\otimes_{\mathscr O_{S_{\mathscr J}}}
  \mathscr O_{Y_{\mathscr J}}
 \right).
\]

\item[(ii)] \textit{The set \(E\) is nonempty if and only if both of the
following conditions hold:}
\begin{enumerate}
\item[(a)] \textit{The problem has a solution locally on \(X\).}
\item[(b)] \textit{A certain obstruction vanishes; it lies in}
\[
 H^1\!\left(
  X_{\mathscr J},
  \operatorname{Hom}_{\mathscr O_{X_{\mathscr J}}}
  \left(
   \mathscr I_{Y_{\mathscr J}},
   \mathscr J\otimes_{\mathscr O_{S_{\mathscr J}}}
   \mathscr O_{Y_{\mathscr J}}
  \right)
 \right).
\]
\end{enumerate}
\end{enumerate}

\medskip
\noindent\SGAInlineRefTarget{sga3:iii:complement:4-4-1}\textbf{Complement 4.4.1.}
\footnote{Editor's note: This complement has been added; it is useful in
proving Proposition 4.8(ii).}
Retain the notation of \SGARefLink{sga3:iii:corollary:4-4}{4.4}, and suppose that \(E\) contains an element
\(Y\). Let \(\mathscr I_Y\) be the ideal of \(\mathscr O_X\) defining
\(Y\), and let \(\mathscr I_{Y_{\mathscr J}}\) be its image in
\(\mathscr O_{X_{\mathscr J}}\). As we saw in the proof of \SGARefLink{sga3:iii:proposition:4-2}{4.2}, there is
a commutative diagram

\begin{center}
  \SGAThreeNativeInput[0.44\textwidth]
    {native_diagrams/exp3/diagram_036_exp3_p046_complement441_square.tex}
\end{center}

\noindent and hence an epimorphism of \(\mathscr O_X\)-modules
\[
 \phi:\mathscr J\mathscr O_X\longrightarrow
 \mathscr J\otimes_{\mathscr O_{S_{\mathscr J}}}
 \mathscr O_{Y_{\mathscr J}}.
\]
Let \(\mathscr K\) be its kernel. For every element \(Y'\) of \(E\),
the morphism \(\mathscr O_X\to\mathscr O_{Y'}\) then factors through
\(\mathscr O_X/\mathscr K\), which is the amalgamated sum \(B'\) in the
proof of \SGARefLink{sga3:iii:lemma:4-1}{Lemma 4.1}. If \(\mathscr I_{Y'}\) denotes the ideal of \(Y'\)
in \(\mathscr O_X\), there is a commutative diagram

\begin{center}
  \SGAThreeNativeInput[0.80\textwidth]
    {native_diagrams/exp3/diagram_037_exp3_p046_complement441_exact_diagram.tex}
\end{center}

Replacing \(X\) by the closed subscheme defined by \(\mathscr K\), we
are thus reduced to the case \(\mathscr K=0\). The datum of \(Y'\) is
then equivalent to that of the \(\mathscr O_X\)-submodule
\(\mathscr I_{Y'}\) of \(\mathscr O_X\) which maps bijectively onto
\(\mathscr I_{Y_{\mathscr J}}\) under the projection
\(p:\mathscr O_X\to\mathscr O_{X_{\mathscr J}}\). Let
\[
 f':\mathscr I_{Y_{\mathscr J}}\xrightarrow{\sim}\mathscr I_{Y'},
 \qquad
 f:\mathscr I_{Y_{\mathscr J}}\xrightarrow{\sim}\mathscr I_Y
\]
denote the respective inverse isomorphisms. Then \(f'-f\) is an element
of
\[
 \operatorname{Hom}_{\mathscr O_{X_{\mathscr J}}}
 \left(
  \mathscr I_{Y_{\mathscr J}},
  \mathscr J\otimes_{\mathscr O_{S_{\mathscr J}}}
  \mathscr O_{Y_{\mathscr J}}
 \right)
 =\operatorname{Hom}_{\mathscr O_{X_{\mathscr J}}}
  (\mathscr I_{Y_{\mathscr J}},\mathscr J\mathscr O_Y),
\]
which we denote by \(d(Y',Y)\). Notice that
\(d(Y,Y')=-d(Y',Y)\).

For our fixed \(Y\) and variable \(Y'\), consider the morphism
\[
 \mathscr I_{Y'}\longrightarrow\mathscr O_X\longrightarrow
 \mathscr O_Y=\mathscr O_X/\mathscr I_Y.
\]
Since its composite with \(\mathscr O_Y\to\mathscr O_{Y_{\mathscr J}}\)
is zero, it takes values in
\[
 \mathscr J\mathscr O_Y
 =\mathscr J\otimes_{\mathscr O_{S_{\mathscr J}}}
  \mathscr O_{Y_{\mathscr J}}.
\]
More precisely, let \(V\) be an open subset of \(X\), let \(x'\) be a
section of \(\mathscr I_{Y'}\) on \(V\), and let \(x_{\mathscr J}\) be
its image in \(\Gamma(V,\mathscr I_{Y_{\mathscr J}})\). Then
\[
 x'=f'(x_{\mathscr J})
   =f(x_{\mathscr J})+(f'-f)(x_{\mathscr J})
   =f(x_{\mathscr J})+d(Y',Y)(x_{\mathscr J}).
\]
Consequently, the morphism
\(\mathscr I_{Y'}\to\mathscr J\mathscr O_Y\) is given by
\(d(Y',Y)\).

